\section{Executable V33 policy specification}
\label{app:policy}

This appendix transcribes the executed notebook rather than reconstructing the
policy from version names. It is pseudocode-equivalent to the six V33 branches;
the archived notebook remains authoritative for field names and floating-point
rounding.

\paragraph{Shared notation and state.}
$\operatorname{clip}(z,a,b)=\max(a,\min(b,z))$ and
$\sigma(z)=(1+e^{-z})^{-1}$. For known horizon $R$ at round $r$,
$t=(r-1)/\max(1,R-1)$; otherwise $t=\min(0.65,(r-1)/16)$. A stall count is the
minimum length of the two players' trailing approximately repeated offer
sequences (relative tolerance $0.003$ in bargaining and $0.001$ in
negotiation). Robust trend is the median adjacent step among the five most
recent observations; volatility is their median absolute deviation.

Memory persists across games only when an opponent name is disclosed. Otherwise
the game ID makes it game-local. Bargaining keys additionally include role,
information regime, own discount, and visible opponent discount; negotiation
keys include role and information regime; persuasion keys include buyer/seller
perspective and signal mode. An event set prevents reprocessing history. Thus
the action is a deterministic function of visible history, game ID, and this
internal memory.

\subsection{Bargaining branches}

For own player $i$, let $d_j$ be the opponent shares in their observed offers,
$\tilde d$ the predicted next demand (last demand plus robust trend), $s$ the
stall count, and $\delta_i$ own discount. The configuration prior $b_0$ is the
Rubinstein responder share computed from visible discounts under complete
information \citep{rubinstein1982perfect};
otherwise it is $0.46+0.02t$ against a disclosed human and
$0.43+0.02t$ otherwise. If $z$ is the largest opponent share in a rejected own
proposal and $m$ the median of the five latest demands, the base floor used
below is $b=\operatorname{clip}(\max\{b_0,z+.006,m-.065+.020t\},.20,.999)$,
omitting unavailable terms.

\paragraph{Alice.}
If demands exist, set
\[
 f=\operatorname{clip}\{\max[f_{\rm reject},
 .75b+.25(\tilde d-.035-.020t)],.18,.90\};
\]
otherwise $f=\operatorname{clip}(b,.18,.90)$. On an offer turn:
\begin{enumerate}
 \item If $s\ge3$, give the responder their last demanded share clipped to
 $[.20,.82]$ (or $[.20,.94]$ on the final round).
 \item Otherwise set $f_s=\operatorname{clip}(f-.010,.18,.88)$,
 $w=\operatorname{clip}(.016+.65\operatorname{MAD}(d),.014,.045)$, and
 $c=.045+.25t+.50(1-\delta_i)$. Search
 $y=0.10,0.105,\ldots,0.94$ and maximize Eq.~\ref{eq:bargaining-search}
 with $\gamma=1.18$.
 \item Under complete information only, transfer an additional
 $\min\{M(.018+.012t),\max(0,x_i-.22M)\}$ from Alice to Bob.
\end{enumerate}
On a response turn, accept a strictly positive final-round offer. Before the
final round, accept if $s\ge3$ and own share is at least $.18$; otherwise accept
when own share is at least
\[
 \max\{.18,.36-.08t-.055s,
 \delta_i(1-f)\operatorname{clip}(.80+g+.10t-.06s,.48,.92)\},
\]
where $g=\operatorname{clip}(.25\times\text{favorable opponent concession},
-.02,.02)$. Under complete information, a rejection is finally overridden
when own share is at least $\max(.20,.33-.10t)$.

\paragraph{Bob.}
Bob uses the same structure with the following replacements:
$f=\operatorname{clip}(\max[f_{\rm reject},.78b+.22(\tilde d-.035-.020t)],.18,.90)$,
$f_s=\operatorname{clip}(f-.006,.18,.88)$,
$w=\operatorname{clip}(.016+.60\operatorname{MAD}(d),.014,.044)$,
$c=.045+.24t+.50(1-\delta_i)$, $\gamma=1.16$, and pre-override risk floor
$\max(.18,.35-.08t-.055s)$. V33 then applies the offer transfer in
Eq.~\ref{eq:bob-transfer}. If the base branch rejects, it accepts whenever own
share is at least
\[
 \max\{.16,(.30\text{ if }\delta_i<.95\text{ else }.33)-.09t-.055s\}.
\]
All proposals divide $M$ exactly. Fixed messages describe the selected split
but do not change it.

\subsection{Bilateral negotiation branches}

Let role $\rho\in\{\text{buyer},\text{seller}\}$, own value $v_i$, latest
opponent price $p_o$, robust median opponent step $\Delta$, and feasible surplus
$S=v_b-v_s$ when both values are visible. Seller utility is $P-v_s$ and buyer
utility is $v_b-P$; these are definitions, not an invocation of the Nash
bargaining solution.

\paragraph{Buyer target.}
Under complete information, start with capture
$a=.66-.10t$, add $.35\operatorname{clip}(-\Delta/S,-.10,.10)$, subtract
$.025\min(s,2)(.35+t)$, and clip $a$ to $[.54,.80]$. Target
$P^*=v_b-aS$. Under hidden information with an opponent offer, project the
seller's price as $\hat p=\max(0,p_o+.5\min(0,\Delta))$ and set
$P^*=v_i-(.64-.09t)(v_i-\hat p)$ when $\hat p\le v_i$, otherwise $P^*=v_i$.
With no offer, $P^*=v_i(.82+.09t)$.

\paragraph{Seller target.}
Under complete information, start with $a=.72-.12t$, subtract $.04$ when
messages are enabled and $.02$ when the horizon is known, add
$.35\operatorname{clip}(\Delta/S,-.10,.10)$, subtract the same stall term, and
clip to $[.54,.80]$. Target $P^*=v_s+aS$. Under hidden information, project a
buyer's price as $\hat p=\max(0,p_o+.5\max(0,\Delta))$ and set
$P^*=v_i+a(\hat p-v_i)$ for $a=.70-.11t-.04\mathbf1_{\rm messages}$ when
$\hat p\ge v_i$, otherwise $P^*=v_i$. With no offer,
$P^*=v_i(1.32-.14t)$.

\paragraph{Responses.}
Both branches accept a final offer iff its utility is nonnegative. The buyer
accepts a profitable repeated path at $s\ge3$; the seller does so at $s\ge2$.
At $s\ge2$, a negative-utility path under an unknown horizon causes walkaway.
Otherwise compare offered utility $U$ with $|P^*-v_i|k$, where
\[
 k_{\rm buyer}=\operatorname{clip}(.76-.13t-.08\min(s,2),.38,.80),
\]
\[
 k_{\rm seller}=\operatorname{clip}(.80-.06\mathbf1_{\rm messages}
 -.13t-.08\min(s,2),.38,.80).
\]
Accept if $U$ exceeds that continuation value or, when $S>0$, captures at least
$.52+.06(1-t)$ (buyer and no-message seller) or $.48+.06(1-t)$ (message
seller). Otherwise counter with $(1-b)P^*+bP$, clipped to the individually
rational side, using
$b=\operatorname{clip}(.22+.38t+.08\min(s,2),0,.70)$ for buyer and adding $.10$
with an $.80$ cap for a message-enabled seller.

For complete-information seller contexts with recorded horizon below five
rounds (including a missing horizon encoded as zero), V33 additionally blends
offers $15\%$ toward the surplus midpoint, accepts nonnegative offers capturing
at least $.40+.08(1-t)$, and blends rejected counters $18\%$ toward the received
price. These historical constants are not estimated equilibrium parameters.

\subsection{Persuasion branches}

\paragraph{Buyer.}
Let high/low values be $v>u$, price $P$, and high-quality prior $p$. Buy if
$P\le u$ and pass if $P>v$. Otherwise parse the seller signal as positive,
negative, or unknown. For a positive signal use $q_H=.90,q_L=.24$; for a
negative signal use $.10,.76$, and first compute
$\pi_0=pq_H/[pq_H+(1-p)q_L]$. Purchased products update high/low counts; passed
products do not update precision because quality is hidden. With
$\alpha=4\pi_0+n_H$ and $\beta=4(1-\pi_0)+n_L$, use the conservative value
\[
 \hat p=\operatorname{clip}\left(
 \frac{\alpha}{\alpha+\beta}-.55
 \sqrt{\frac{\alpha\beta}{(\alpha+\beta)^2(\alpha+\beta+1)}},0,1\right).
\]
For an unknown signal use $p$. Buy when
$\hat p v+(1-\hat p)u+B\ge P$, where $B=.010(v-u)$ times the remaining-round
fraction divided by $\sqrt{1+n}$ only for a positive signal within $.08(v-u)$
of the cutoff and with fewer than four observed labels; otherwise $B=0$.

\paragraph{Seller.}
Always signal positive for high quality. For low quality, compute the target in
Eq.~\ref{eq:static-calibration}; return $q=1$ when $c\le0$, $q=0$ when $c\ge1$,
and $q=0$ for an interior cutoff with $p\in\{0,1\}$. When buyer values are
hidden use the ad hoc rule reported in the main text. The notebook defines
terminal as $r\ge R$ after replacing a missing $R$ with the current round; thus
an absent horizon receives the terminal $.05$ increment and $.94$ utilization.
If $n$ low-quality products precede the current one, signal positive exactly
when $\lfloor(n+1)q\rfloor>\lfloor nq\rfloor$. Text mode emits fixed BUY/PASS
sentences and binary mode emits yes/no.

For constant $q$, the floor-difference sequence has prefix discrepancy below
one. V33 recomputes $q$ and reveals signals, so this arithmetic fact supplies no
dynamic persuasion or incentive-compatibility guarantee.

\subsection{Contract wrapper and exposure controller}

The wrapper rejects unexpected keys, nonfinite or negative amounts, allocations
that do not sum to the pot within tolerance, illegal decisions, missing
nonterminal counteroffers, and messages over 2,000 characters. Its fallbacks
are: equal split/accept in bargaining; own-value offer, profitable acceptance,
or own-value counter in negotiation; PASS/no as seller; and prior expected-value
purchase as buyer. Every fallback is logged.

The controller starts only at zero active games, queues one family, leaves the
queue after observing an assignment, drains it, and checks that assignment and
authoritative count increments are at most one. Family baseline-loss limits for
bargaining/negotiation/persuasion are $30/40/28$ and peak drawdowns are
$24/30/22$ after 20 games. Role limits are $24/28/20$ and role drawdowns
$18/22/16$ after ten role games. Three consecutive matchmaking timeouts stop a
family. The executed notebook measured new telemetry before reward harvesting;
the 500 caught HTTP 422 errors appended during harvesting therefore escaped the
abort check. This historical monitoring-order bug is preserved in the executed
artifact and corrected only in the documented future protocol.

\paragraph{Version map.}
``Alice V27'' denotes the Alice branch plus its complete-information correction;
``Bob V29'' denotes the Bob branch plus delay-aware transfer/acceptance override;
``buyer V23 rollback'' denotes the negotiation buyer equations above;
``seller V27'' denotes the seller branch plus short-horizon repair; and
``persuasion buyer V23'' denotes the purchased-only conservative posterior rule.
These labels record provenance and are not separate algorithms required to
reconstruct V33.
